% Este fuente se publica en dos dependencias causales mediante los
% selectores definidos por los wrappers de integración. Ambos tramos cubren
% conjuntamente el archivo completo, sin solape ni omisión.
\ifdefined\HMTGMCOperatorialAPP
\section[Image matricielle d'APP]{Image matricielle d'APP : classification exacte
de l'algèbre somme--produit et intervalle libre}
\label{p12-83-c22-s1:gmc:chap:operatorios}

Les carrés qutrit du \cref{p12-83-c22-s3:gmc:chap:qms} fournissent une entrée
opératorielle. Il existe une
seconde entrée, plus interne à APP : les deux partitions induites par la somme et le
produit engendrent des projections orthogonales non commutatives. Le langage des
systèmes d'opérateurs permet de les étudier sans imposer une grille.

\subsection{Couple d'espérances conditionnelles associé aux partitions additive et multiplicative d'APP}

Soit
\[
 \mathcal H_d=\ell^2(\{1,\ldots,9\}^d),
\]
et définissons
\[
 \Sigma_d(x)=\operatorname{dr}(x_1+\cdots+x_d),
 \qquad
 \Pi_d(x)=\operatorname{dr}(x_1\cdots x_d).
 \tag{GMC-6.1}
\]
Soient $P_{\Sigma,d}$ et $P_{\Pi,d}$ les projections orthogonales sur les
fonctions constantes sur les fibres de $\Sigma_d$ et $\Pi_d$. Le pont
précédent démontre
\[
 \|[P_{\Sigma,2},P_{\Pi,2}]\|=\frac{\sqrt2}{3},
 \qquad
 \|[P_{\Sigma,3},P_{\Pi,3}]\|=\frac{\sqrt3}{4},
 \tag{GMC-6.2}
\]
et
\[
 \operatorname{Ran}P_{\Sigma,d}\cap
 \operatorname{Ran}P_{\Pi,d}=\mathbb C\mathbf1
 \qquad(2\le d\le15).
 \tag{GMC-6.3}
\]
APP contient donc un ensemble opératoriel non commutatif avant
toute dénomination QLS ou QMS.

On définit
\[
 \mathcal S_d=\Span\{I,P_{\Sigma,d},P_{\Pi,d}\},
 \qquad
 \mathfrak C_d^{\rm APP}=C^*(P_{\Sigma,d},P_{\Pi,d}).
 \tag{GMC-6.4}
\]
$\mathcal S_d$ est un système d'opérateurs et son image matricielle est
\[
 \mathcal W_s(\mathrm{APP}_d)=
 \left\{(\Phi(P_{\Sigma,d}),\Phi(P_{\Pi,d})):
 \Phi:\mathfrak C_d^{\rm APP}\to M_s\ \text{UCP}\right\}.
 \tag{GMC-6.5}
\]
La tour $\bigsqcup_s\mathcal W_s$ est matriciellement convexe. À la différence d'une
étiquette générale de « quanticité », \textup{(GMC-6.5)} permet de calculer des extrêmes,
des compressions, des dilatations et des invariants.

\subsection{Classification de l'algèbre somme--produit à la profondeur 2}

Posons $P=P_{\Sigma,2}$ et $Q=P_{\Pi,2}$ sur $\mathcal H_2$, de dimension
81. Les neuf fibres de $\Sigma_2$ sont de taille neuf. Les fibres de
$\Pi_2$ ont pour tailles
\[
 6,6,12,6,6,12,6,6,21.
 \tag{GMC-6.6}
\]
Soient $u_a$ et $v_b$ les vecteurs indicateurs normalisés de ces fibres, et
$C=(\langle u_a,v_b\rangle)_{a,b=1}^9$ leur matrice de Gram croisée. Le dénombrement
exact des intersections produit
\[
 \chi_{CC^*}(t)
 =t^5(t-1)\left(t-\frac57\right)
       \left(t-\frac13\right)^2.
 \tag{GMC-6.7}
\]
De plus,
\[
\begin{array}{c|c}
\text{sous-espace commun} & \text{dimension}\\
\hline
\ker P\cap\ker Q&64\\
\operatorname{Ran}P\cap\operatorname{Ran}Q&1\\
\operatorname{Ran}P\cap\ker Q&5\\
\ker P\cap\operatorname{Ran}Q&5
\end{array}
\tag{GMC-6.8}
\]

\begin{theorem}[Algèbre somme--produit d'APP dans $d=2$]
\label{p12-83-c22-s1:gmc:thm:algebra-app-d2}
On a un isomorphisme de $C^*$-algèbres de dimension finie
\[
 \boxed{\ 
 \mathfrak C_2^{\rm APP}
 \cong\mathbb C^4\oplus M_2(\mathbb C)\oplus M_2(\mathbb C).
 \ }
 \tag{GMC-6.9}
\]
En particulier, $\dim\mathfrak C_2^{\rm APP}=12$.
\end{theorem}

\begin{proof}
La décomposition canonique de Halmos pour deux projections sépare les quatre
sous-espaces de \textup{(GMC-6.8)} et un secteur générique pour chaque angle principal
$0<\lambda<1$. L'équation \textup{(GMC-6.7)} montre que les seules valeurs
génériques distinctes sont $5/7$ et $1/3$. Dans chacun d'eux, le couple est unitairement
équivalent à
\[
 P_0=\begin{pmatrix}1&0\\0&0\end{pmatrix},
 \qquad
 Q_\lambda=
 \begin{pmatrix}
  \lambda&\sqrt{\lambda(1-\lambda)}\\
  \sqrt{\lambda(1-\lambda)}&1-\lambda
 \end{pmatrix},
 \tag{GMC-6.10}
\]
qui engendre $M_2$. La multiplicité de $\lambda=1/3$ amplifie la même
représentation, mais ne crée pas d'autre terme simple. Les quatre secteurs
simultanément $0/1$ produisent $\mathbb C^4$.
\end{proof}

Plus précisément,
\[
 Q_{5/7}=
 \begin{pmatrix}5/7&\sqrt{10}/7\\\sqrt{10}/7&2/7\end{pmatrix},
 \qquad
 Q_{1/3}=
 \begin{pmatrix}1/3&\sqrt2/3\\\sqrt2/3&2/3\end{pmatrix}.
 \tag{GMC-6.11}
\]
La norme du commutateur est le maximum de
$\sqrt{\lambda(1-\lambda)}$ et retrouve $\sqrt2/3$.

\begin{corollary}[Paramétrage complet de l'image matricielle]
\label{p12-83-c22-s1:gmc:cor:rango-app-d2}
Un couple $(X,Y)$ appartient à $\mathcal W_s(\mathrm{APP}_2)$ si et seulement s'il
existe des effets
$E_{00},E_{11},E_{10},E_{01}\succeq0$ et des applications CP
$\Psi_{5/7},\Psi_{1/3}:M_2\to M_s$ tels que
\[
 E_{00}+E_{11}+E_{10}+E_{01}
 +\Psi_{5/7}(I)+\Psi_{1/3}(I)=I_s,
 \tag{GMC-6.12}
\]
\[
 X=E_{11}+E_{10}
 +\Psi_{5/7}(P_0)+\Psi_{1/3}(P_0),
 \tag{GMC-6.13}
\]
\[
 Y=E_{11}+E_{01}
 +\Psi_{5/7}(Q_{5/7})+\Psi_{1/3}(Q_{1/3}).
 \tag{GMC-6.14}
\]
\end{corollary}

\begin{proof}
Toute application CP issue d'une somme directe se décompose en applications CP sur ses
termes. Les quatre termes scalaires produisent les effets $E_{ab}$ ; les
deux termes matriciels produisent $\Psi_\lambda$. L'unitalité est
\textup{(GMC-6.12)} ; les coordonnées de $P$ et $Q$ dans \textup{(GMC-6.9)} donnent
\textup{(GMC-6.13)}--\textup{(GMC-6.14)}. L'argument est réversible.
\end{proof}

\begin{bridgebox}
La classification \textup{(GMC-6.9)} fournit la réalisation opératorielle du pont. Elle transforme
l'incompatibilité somme--produit en un objet opératoriel entièrement
calculable : quatre secteurs classiques et deux blocs matriciels irréductibles.
La valeur $1/3$ réapparaît ici comme angle principal opératoriel, non seulement comme
recouvrement MUB ; $5/7$ est la seconde valeur générique de $PQP$ dans le secteur de
somme--produit. Aucune ne doit être identifiée à la contraction $5/9$ du bloc
central $B_c$. Ce sont des invariants d'objets distincts, coordonnés sans
être identifiés par ressemblance décimale.
\end{bridgebox}

\subsection{Séparation spectrale \texorpdfstring{$4$}{4} et intervalle matriciel libre}

Pour le bloc de \textup{(GMC-2.4)},
$B_c=7I+2R$ avec $R=R^*$ non scalaire, $R^2=I$ et
$\operatorname{Spec}(R)=\{-1,+1\}$,
$C^*(B_c)=C^*(R)\cong\mathbb C^2$. Son image matricielle au niveau $s$ est
\[
 \mathcal W_s(B_c)=
 \{\Phi(B_c):\Phi:C^*(B_c)\to M_s\text{ UCP}\}.
\]

\begin{theorem}[Intervalle libre du bloc central]
\label{p12-83-c22-s1:gmc:thm:intervalo-bc}
Pour tout $s\ge1$,
\[
 \boxed{\ 
 \mathcal W_s(B_c)
 =\{Y=Y^*:5I_s\preceq Y\preceq9I_s\}.
 \ }
 \tag{GMC-6.15}
\]
\end{theorem}

\begin{proof}
Une application UCP envoie la symétrie $R$ sur une contraction autoadjointe
$-I_s\preceq X\preceq I_s$ et, par conséquent,
$\Phi(B_c)=7I_s+2X$ appartient à l'intervalle. Réciproquement, étant donné un tel $X$,
les effets $(I_s\pm X)/2$ définissent une application UCP issue de $\mathbb C^2$ qui envoie
$R$ sur $X$.
\end{proof}

Ainsi, $7$ est le centre, $2$ le rayon opératoriel et $4$ le diamètre spectral
de l'intervalle libre. Le bord positif, le bord de l'intervalle matriciel et
l'hyperbole de déterminant un de \textup{(GMC-2.6)} sont trois notions distinctes qui
peuvent désormais être coordonnées sans être confondues.

La double projection ne réside pas dans cette classification précoce. Sa formulation
à partir d'un antécédent spectral commun est publiée, avec ses hypothèses et ses limites,
dans la section dédiée du chapitre directeur de la double projection.
\fi

\ifdefined\HMTGMCDoubleProjection
\subsection{Antécédent commun des lecteurs : raffinement spectral et limite de
dilatation conjointe}

Soit $X_m$ un ensemble fini d'histoires et soient
\[
 R_m:X_m\to\mathcal A_m,
 \qquad
 Q_m:X_m\to\mathcal E_m
 \tag{GMC-6.16}
\]
deux lecteurs : par exemple, l'un archimédien et l'autre exceptionnel. Chacun
définit une \PVM dans $C(X_m)$ au moyen des fonctions caractéristiques de ses
fibres.

\begin{proposition}[Antécédent spectral commun]
\label{p12-83-c22-s1:gmc:prop:doble-proyeccion-conmutativa}
Les deux \PVM de \textup{(GMC-6.16)} commutent et possèdent le raffinement commun
\[
 P_{a,e}=1_{\{R_m=a\}}1_{\{Q_m=e\}}
 =1_{\{x:R_m(x)=a,\ Q_m(x)=e\}}.
 \tag{GMC-6.17}
\]
Si les lecteurs sont compatibles avec les troncatures, les raffinements sont
compatibles dans $m$ et déterminent deux lecteurs sur $C(X_\infty)$ avec une
\PVM cylindrique commune.
\end{proposition}

\begin{proof}
Toutes les fonctions appartiennent à la même algèbre commutative. Le produit
des indicatrices est l'indicatrice de l'intersection. Sommer sur $e$ retrouve
la \PVM de $R_m$ et sommer sur $a$ retrouve celle de $Q_m$. La compatibilité passe
à la limite par le \Cref{p12-83-c22-s4:gmc:thm:dilatacion-diagonal}.
\end{proof}

Ce résultat formalise l'antécédent commun de la double projection. Il ne
démontre pas que le couple de lecteurs soit conjointement injectif : pour cela, il doit
séparer les points de $X_\infty$. Il ne construit pas non plus automatiquement un antécédent
non commutatif commun.

\begin{proposition}[Incompatibilité de deux mesures projectives mutuellement non biaisées]
Deux \PVM qutrit mutuellement non biaisées n'admettent pas de mesure conjointe nette
dans $M_3(\mathbb C)$.
\end{proposition}

\begin{proof}
Deux \PVM finies sont conjointement mesurables par une \PVM commune si et seulement si
elles commutent. Des projecteurs de contextes MUB distincts ont un recouvrement $1/3$ et
ne commutent pas.
\end{proof}


La double projection HMT ne doit donc pas être identifiée à deux contextes
MUB nets sur le même qutrit. Elle peut se réaliser dans la diagonale commune des
histoires, au moyen d'effets bruités, ou dans un antécédent non commutatif plus
vaste avec deux canaux UCP. La dernière possibilité exige l'inclusion des images
matricielles, la positivité de Choi et l'équivariance par rapport à la dynamique TPK.

\begin{workbox}[fontupper=\footnotesize,fonttitle=\bfseries\footnotesize,
  boxsep=1pt,left=4pt,right=4pt,top=0pt,bottom=0pt,
  toptitle=1pt,bottomtitle=1pt]
L'antécédent non commutatif est construit dans
\Cref{sec:cpe-antecedente-ucp} : le système $\mathcal S_\infty$ et les applications UCP
$\Phi_{\rm arq}$ et $\Phi_{\rm exc}$ reproduisent simultanément les deux
projections HMT et entrelacent la dynamique TPK dans le domaine équivariant
déclaré. La \Cref{p12-83-c22-s1:gmc:prop:doble-proyeccion-conmutativa}
fournit son antécédent diagonal.
\end{workbox}
\fi


\section{Convexité matricielle généalogique : objets, domaines et dictionnaire
typé}

\subsection{Sous-structures, objet global et projections}

La question qui organise le travail de De las Cuevas--Netzer est de savoir comment une
collection de mesures locales peut composer un objet bidimensionnel et jusqu'à
quel point cet objet est la compression d'une totalité idéale. Une mesure
isolée ---une \POVM--- admet une dilatation de Naimark ou de Stinespring. Toutefois,
une grille contient simultanément une \POVM dans chaque ligne et dans
chaque colonne. Demander une dilatation commune qui rende projectives toutes ces
mesures est une condition globale beaucoup plus forte. Il existe des carrés magiques
quantiques qui ne la satisfont pas, même au plus petit niveau matriciel non commutatif
pertinent \citep{delascuevas2020qms}.

HMT formule un autre problème de parties et de totalités. La coordonnée visible d'un
état ne conserve pas nécessairement la route qui l'a produite. APP distingue les
évaluations somme et produit sur un support commun ; TRIT conserve orientation
et retenue ; TPK conserve composition des chemins, retour, holonomie et mémoire.
Le continu naît du prolongement cohérent des histoires finies, et la
double projection publie deux ombres différentes d'un même état
enrichi. Les deux théories demandent donc ce qui se perd lors de la compression.
La différence est que l'une organise la perte par niveau matriciel et l'autre par
profondeur généalogique.

\begin{intuitionbox}
De las Cuevas demande : « puis-je voir cette grille imparfaite comme la
compression d'une grille projective plus grande ? ». HMT demande : « quelle histoire
complète a été comprimée pour produire cette sortie ? ». Le pont répond que
les deux questions peuvent être formulées sur la même famille, avec deux indices
indépendants : profondeur $m$ et dimension d'observation $s$.
\end{intuitionbox}

\subsection{Carrés quantiques, systèmes d'opérateurs et convexité matricielle}

Soient $n,s\ge1$. Un carré magique quantique de taille extérieure $n$ et
intérieure $s$ est une matrice
\[
 A=(A_{ij})_{i,j=0}^{n-1}\in\Mat_n(\Her_s(\mathbb C))
\]
tel que
\[
 A_{ij}\succeq0,\qquad
 \sum_{j=0}^{n-1}A_{ij}=I_s,\qquad
 \sum_{i=0}^{n-1}A_{ij}=I_s.
 \tag{GMC-1.1}
\]
Chaque ligne et chaque colonne est une \POVM. Si toutes les entrées sont des projections,
$A$ est une matrice de permutation quantique ; si, en outre, ces projections
commutent, $A$ est une matrice de permutation quantique commutative
\citep{delascuevas2020qms,delascuevas2023latin}.

Un \QMS est \emph{semi-classique} s'il existe une \POVM
\(\{Q_\sigma\}_{\sigma\in S_n}\) telle que
\[
 A=\sum_{\sigma\in S_n}P_\sigma\otimes Q_\sigma,
 \tag{GMC-1.2}
\]
où $P_\sigma$ est la matrice de permutation classique associée à
$\sigma$. La formule n'exige pas que les effets $Q_\sigma$ commutent. Elle dit
que toute la grille s'obtient en réordonnant une même mesure primaire.

Un carré latin quantique d'ordre $n$ est une grille de vecteurs
unitaires $v_{ij}\in\mathbb C^n$ dont les lignes et les colonnes sont des bases
orthonormées. Les projecteurs de rang un
\[
 P_{ij}=|v_{ij}\rangle\langle v_{ij}|
\]
forment un \QMS projectif. Les \QLS de classe facile sont ceux produits par un carré
latin classique et une base orthonormée fixe. De las Cuevas, Netzer et
Valentiner-Branth démontrent que les \QLS semi-classiques sont exactement ceux de
classe facile et que l'enveloppe matriciellement convexe des \QLS contient davantage que les
carrés semi-classiques \citep{delascuevas2023latin}.

L'opération matriciellement convexe élémentaire est
\[
 X=\sum_{r=1}^{k}V_r^*X^{(r)}V_r,
 \qquad
 \sum_{r=1}^{k}V_r^*V_r=I_s.
 \tag{GMC-1.3}
\]
Elle permet de combiner des objets de tailles intérieures différentes. En langage
opératoriel, les compressions et les sommes directes s'organisent au moyen
d'applications unitales complètement positives. L'ensemble de toutes les
valeurs matricielles d'un système d'opérateurs est son image matricielle.

\begin{importbox}
Les résultats directeurs utilisés sont les suivants : un \QMS se dilate en une matrice de
permutation quantique à entrées commutatives si et seulement s'il est semi-classique ;
pour $n\ge3$ et $s\ge2$, il existe des \QMS non semi-classiques, et pour tout $n\ge3$, il existe
des \QMS sur $M_2(\mathbb C)$ qui ne se dilatent en aucune matrice de
permutation quantique \citep{delascuevas2020qms}. La dilatabilité simultanée
est donc une propriété substantielle, non une conséquence automatique du
fait que chaque ligne et chaque colonne se dilatent séparément.
\end{importbox}

\subsection{Correspondance typée entre l'état généalogique HMT–MD et les systèmes d'opérateurs : profondeur, troncature, monodromie et lecture}

Dans HMT, une sortie visible n'épuise pas l'état. Une section finie
porte, selon le domaine, des données du type
\[
 x_m=(\text{préfixe},\text{feuillet},\text{orientation},\text{retenue},
       \text{frontière},\text{route},\text{mémoire},\text{jauge}).
 \tag{GMC-1.4}
\]
La profondeur $m$ n'est pas une dimension de Hilbert. Elle compte la part de la
généalogie qui a été résolue. Deux états peuvent partager une coordonnée visible et
différer en $m$, en feuillet ou en mémoire.

APP fournit l'orthogramme nonadique et les feuillets somme et produit. TRIT
enregistre l'orientation ternaire et la retenue locale. TPK est le noyau
topologique de phase et la catégorie de routes qui compose ces données. Dans une
réalisation finie, les routes agissent par des opérateurs sur un espace de
Hilbert, mais la représentation ne remplace pas la généalogie : elle la transporte.

Le langage de HMT contient en outre trois opérations qui ne sont pas présentes dans
la définition nue d'un \QMS :

\begin{enumerate}
  \item \emph{troncature} : oublier la continuation postérieure à la profondeur
  $m$, en conservant le préfixe et en agrégeant ses prolongations ;
  \item \emph{monodromie} : revenir à la même phase visible avec une mémoire non
  triviale ;
  \item \emph{lecture} : publier une coordonnée mathématique ou physique sans
  l'identifier à la section complète.
\end{enumerate}

\Needspace{.74\textheight}
\subsection{Correspondances directes et inverses entre les formalismes}

\begin{longtable}{@{}L{.23\textwidth}L{.29\textwidth}L{.37\textwidth}@{}}
\toprule
Langage matriciel & Lenguaje HMT & Correspondencia tipada\\
\midrule
\endfirsthead
\toprule
Lenguaje matricial & Lenguaje HMT & Correspondencia tipada\\
\midrule
\endhead
Taille extérieure $n$ & grille publiée & nombre de lignes et de colonnes de
l'objet magique ; ce n'est ni la profondeur ni la dimension intérieure\\
Taille intérieure $s$ & carta operatorial & dimension des effets ou de la
compression matricielle\\
Profundidad $m$ & généalogie & nouvel indice qui ne se réduit ni à $n$ ni à
$s$\\
\POVM & famille de lecteurs positifs & effets dont la somme est l'unité ; la somme
exprime la normalisation, non la mémoire complète\\
\PVM & contexto espectral & famille de projections orthogonales qui résout
l'identité\\
\UCP & lecteur matriciel & application qui conserve l'unité et la positivité à tous les
niveaux\\
Compression $V^*(\cdot)V$ & ombre & passage d'une réalisation plus grande à une
carte d'observation plus petite\\
Dilatation & antecedente operatorio & représentation multiplicative commune
avant compression\\
Semi-classicité & une \POVM primaire commune réordonnée & propriété formelle ; elle
n'équivaut pas à la commutativité de toutes les cellules\\
No semiclasicidad & composition véritablement bidimensionnelle & ne peut pas
se réduire au réordonnancement d'une \POVM primaire commune\\
Système d'opérateurs & espace des lecteurs & sous-espace autoadjoint contenant
l'unité qui organise toutes les cartes UCP\\
Rango matricial & ensemble des ombres à tous les $s$ & famille libre
d'évaluations internes\\
Limite inductive & clôture des raffinements opératoires & conservation de
l'unité à travers toute la filtration nonadique\\
Limite projective & section infinie compatible & histoire qui conserve tous
ses préfixes\\
\bottomrule
\end{longtable}

\subsection{Extension de la bigraduation à 3 indices}

Le formalisme des carrés magiques sépare $(n,s)$. HMT ajoute $m$.
L'objet total doit s'écrire
\[
 \mathfrak M^{(n)}_{m,s},
 \tag{GMC-1.5}
\]
et non simplement \(\mathfrak M_n\) ou \(\mathfrak M_s\). Chaque indice répond
à une question différente :

\begin{center}
\begin{tabularx}{.92\textwidth}{@{}c L{.25\textwidth}Y@{}}
\toprule
Indice & Question & Opération associée\\
\midrule
$m$ & quelle quantité d'histoire est conservée ? & raffinement et troncature\\
$n$ & comment la grille est-elle composée ? & lignes, colonnes et permutations\\
$s$ & depuis quel niveau matriciel observe-t-on ? & somme directe et compression UCP\\
\bottomrule
\end{tabularx}
\end{center}

Pour APP, on fixe $n=9$. Alors la famille
$\{\mathfrak M^{(9)}_{m,s}\}_{m,s}$ est bigraduée par $(m,s)$. Si l'on
compare APP, les quatre \QLS d'ordre trois et l'atlas des douze directions,
l'indice $n$ varie de nouveau et la structure correcte est trigraduée.

\begin{controlbox}
Parler de « bigraduation $(m,s)$ » n'efface pas $n$ : cela signifie que l'on étudie une
famille de taille extérieure fixée. Appeler APP « $9\times9$ » ne fixe pas non plus
la taille intérieure de ses opérateurs. Les matrices internes $6\times6$, la
carte qutrit $s=3$ et la grille extérieure $n=9$ sont des objets de types
distincts.
\end{controlbox}

\subsection{Bidirectionnalité de l'échange mathématique}

La direction De las Cuevas--Netzer $\longrightarrow$ HMT apporte des critères que HMT n'avait pas
formulés avec ce degré de séparation : positivité cellule par cellule,
normalisation de toutes les lignes et colonnes, semi-classicité, frontière de
l'enveloppe matriciellement convexe, dilatation simultanée et points extrémaux d'Arveson. Le
résultat est un ensemble de preuves et de réfutateurs pour toute allégation de
« carré quantique ».

La direction HMT $\longrightarrow$ De las Cuevas--Netzer ajoute de la structure à l'objet qui
est comprimé : profondeur, route, non-scindement de la retenue, monodromie, dynamique,
jauge, double projection et lecteurs métrologiques. Le problème ne consiste plus
seulement à savoir s'il existe une dilatation, mais à demander s'il existe une dilatation
\emph{généalogiquement fidèle} : une dilatation qui n'identifie pas deux routes distinctes au seul motif
qu'elles produisent le même effet visible.

\begin{bridgebox}
La notion nouvelle qui organise l'article est la \emph{dilatation simultanée
généalogique} : une représentation multiplicative commune qui dilate non seulement plusieurs
lignes et colonnes à profondeur fixe, mais toutes les profondeurs compatibles,
et qui spécifie quelle partie de la mémoire demeure dans la représentation et quelle
partie se perd dans la compression. Le théorème exact est démontré dans
\Cref{p12-83-c22-s4:gmc:chap:bigraduacion}.
\end{bridgebox}


\section[Systèmes opératoriels d'incidence]{Systèmes opératoriels
d'incidence : carrés latins quantiques, contextes qutrit et dilatation
généalogique commune}
\label{p12-83-c22-s3:gmc:chap:qms}

L'objet extérieur d'APP est d'ordre neuf, mais le nombre neuf ne suffit pas
à en faire un carré magique quantique. Il faut présenter 81 effets
positifs et vérifier 18 normalisations opératorielles. Ce chapitre réalise
cette tâche pour deux ombres d'APP et pour l'atlas complet des douze
directions. Il détermine également leur classe exacte dans la taxonomie de De las
Cuevas--Netzer.

\subsection{4 carrés latins quantiques cycliques de classe opératorielle facile}

Soient
\[
 \mathcal P_a=\{P_{a,0},P_{a,1},P_{a,2}\},
 \qquad a\in\mathbb P^1(\mathbb F_3),
 \tag{GMC-4.1}
\]
les quatre \PVM qutrit du système de Weyl. Pour chaque $a$, choisissez des vecteurs
unitaires $v_{a,k}$ avec
$P_{a,k}=|v_{a,k}\rangle\langle v_{a,k}|$ et définissez
\[
 L^{(a)}_{ij}=v_{a,i+j\bmod3}.
 \tag{GMC-4.2}
\]
Chaque ligne et chaque colonne parcourt une fois la base
$(v_{a,0},v_{a,1},v_{a,2})$. Par conséquent, $L^{(a)}$ est un \QLS d'ordre trois.
Puisqu'il est obtenu à partir de la table latine cyclique et d'une base orthonormée fixée, il est de classe facile et son
\QMS projectif est semi-classique
\citep{delascuevas2023latin}.

\begin{proposition}[Obstruction MUB à l'empilement latin]
\label{p12-83-c22-s3:gmc:prop:no-go-mub}
On ne peut pas prendre deux lignes de contextes distincts
$\mathcal P_a$ et $\mathcal P_b$, $a\ne b$, comme lignes d'un même \QLS
d'ordre trois en conservant leurs colonnes alignées.
\end{proposition}

\begin{proof}
La condition de colonne d'un \QLS exigerait l'orthogonalité des deux
vecteurs d'une colonne. Mais l'absence de biais mutuel donne
\[
 |\langle v_{a,k},v_{b,\ell}\rangle|^2
 =\Tr(P_{a,k}P_{b,\ell})=\frac13
\]
pour tout $k,\ell$. Le recouvrement n'est pas nul.
\end{proof}

L'obstruction ne nie pas la compatibilité quantique des quatre bases.
Elle détermine que la structure opératorielle appropriée pour les utiliser simultanément n'est pas un
\QLS de rang un : il faut passer à des effets pondérés et à un \QMS.

\subsection{Lemme cyclique des carrés magiques quantiques associés à une mesure opératorielle positive}

\begin{lemma}[Carré cyclique associé à une \POVM]
\label{p12-83-c22-s3:gmc:lem:qms-ciclico}
Soit $E=(E_0,\ldots,E_{n-1})$ une \POVM sur $\mathbb C^s$. Alors
\[
 A_{ij}=E_{i+j\bmod n},
 \qquad 0\le i,j<n,
 \tag{GMC-4.3}
\]
est un \QMS de tailles $(n,s)$. Il est en outre semi-classique et admet une
dilatation simultanée en un \QPM commutatif.
\end{lemma}

\begin{proof}
Chaque entrée est positive. Toute ligne et toute colonne est une permutation des
$n$ effets, de sorte que leur somme est $I_s$. Pour $t\in\mathbb Z/n\mathbb Z$,
soient $\sigma_t(i)=t-i$ et $P_{\sigma_t}$ sa matrice de permutation. Alors
\[
 A=\sum_{t=0}^{n-1}P_{\sigma_t}\otimes E_t,
 \tag{GMC-4.4}
\]
ce qui est une décomposition semi-classique explicite. Si
$E_t=V^*Q_tV$ est une dilatation de Naimark de la \POVM en une \PVM
$(Q_t)_t$, les entrées
$\widehat A_{ij}=Q_{i+j}$ forment un \QPM commutatif et
$A_{ij}=V^*\widehat A_{ij}V$ simultanément pour les $n^2$ cellules.
\end{proof}

Cette construction appartient à la classe des familles latines associées à des POVM
(\emph{POVM Latin-square friends})
de De las Cuevas et Netzer \citep{delascuevas2023latin}. La nouveauté du pont n'est pas le lemme
abstrait, mais les \POVM que HMT fournit, leur étalonnage nonadique et la
séparation exacte entre ce que chaque carré conserve et ce qu'il
comprime encore.

\subsection{Représentation additive d'APP de tailles \((9,3)\) et son codomaine d'incidence}

Fixons l'un des quatre contextes, par exemple
$\mathcal P_0=\{P_{0,0},P_{0,1},P_{0,2}\}$. Pour $t\in\mathbb Z/9\mathbb Z$,
nous définissons
\[
 E_t^{\Sigma}=\frac13P_{0,t\bmod3},
 \qquad
 A_{ij}^{\Sigma}=E_{i+j\bmod9}^{\Sigma}.
 \tag{GMC-4.5}
\]
Chaque projection apparaît trois fois, donc
\[
 \sum_{t=0}^{8}E_t^{\Sigma}
 =\frac13\sum_{k=0}^2 3P_{0,k}=I_3.
\]
D'après le \Cref{p12-83-c22-s3:gmc:lem:qms-ciclico}, $A^\Sigma$ est un \QMS semi-classique de tailles
$(9,3)$. C'est une transduction directe de la loi latine
$\Sigma(i,j)=i+j\bmod9$ vers une lecture qutrit. Il conserve la taille extérieure et
le feuillet additif ; délibérément, il ne distingue pas les trois relèvements de Hensel d'une
même classe ternaire.

\subsection{\(3\) contextes dans un carré $(9,3)$}

Choisissons trois des quatre \PVM qutrit, indexées par $a=0,1,2$, et ordonnons
leurs projecteurs au moyen de $t=3a+k$ :
\[
 E_{3a+k}^{(9)}=\frac13P_{a,k},
 \qquad a,k\in\mathbb Z/3\mathbb Z.
 \tag{GMC-4.6}
\]
Comme
\[
 \sum_{a,k}E_{3a+k}^{(9)}
 =\frac13\sum_{a=0}^2\sum_{k=0}^2P_{a,k}=I_3,
\]
la formule
\[
 A_{ij}^{(9)}=E_{i+j\bmod9}^{(9)}
 \tag{GMC-4.7}
\]
produit un autre \QMS de tailles $(9,3)$.

\begin{theorem}[Carré qutrit d'ordre extérieur neuf]
\label{p12-83-c22-s3:gmc:thm:qms-9-3}
Le carré $A^{(9)}$ de \textup{(GMC-4.7)} satisfait :

\begin{enumerate}
 \item ses 81 cellules sont des effets positifs de rang un et de trace $1/3$ ;
 \item toutes les lignes et colonnes ont pour somme $I_3$ ;
 \item il contient des entrées qui ne commutent pas ;
 \item il est semi-classique au moyen de
 \[
 A^{(9)}=\sum_{t=0}^{8}P_{\sigma_t}\otimes E_t^{(9)};
 \tag{GMC-4.8}
 \]
 \item il admet une dilatation simultanée commutative de ses 81
 entrées.
\end{enumerate}
\end{theorem}

\begin{proof}
Les deux premiers points et la décomposition découlent du
\Cref{p12-83-c22-s3:gmc:lem:qms-ciclico}. Si $a\ne b$, deux projecteurs MUB de rang un ne
commutent qu'en cas de coïncidence ou d'orthogonalité, impossibles puisque leur recouvrement
est $1/3$ ; certains
$[E_{3a+k}^{(9)},E_{3b+\ell}^{(9)}]$ sont donc non nuls. La dernière affirmation est la
dilatation construite dans la preuve du lemme.
\end{proof}

\begin{bridgebox}
Ce théorème clôt une existence que le pont préliminaire avait laissée
ouverte : on dispose désormais de 81 effets explicites avec $n=9$, $s=3$ et des normalisations
exactes. Il démontre en même temps une séparation conceptuelle féconde :
\[
 \text{cellules non commutatives}
 \quad\centernot\Longrightarrow\quad
 \text{non-semi-classicité}.
\]
Ici, les cellules ne commutent pas, mais toute la grille provient du réordonnancement d'une
seule \POVM primaire.
\end{bridgebox}

\subsection{L'atlas complet dans un carré \texorpdfstring{$(12,3)$}{(12,3)}}

La droite projective sur l'anneau nonadique possède douze points :
\[
 |\mathbb P^1(\mathbb Z/9\mathbb Z)|=12.
 \tag{GMC-4.9}
\]
La réduction modulo trois induit
\[
 r:\mathbb P^1(\mathbb Z/9\mathbb Z)
 \longrightarrow\mathbb P^1(\mathbb F_3),
 \tag{GMC-4.10}
\]
avec quatre fibres de cardinal trois. Une fois fixé un étalonnage
\[
 \Xi:\mathbb P^1(\mathbb Z/9\mathbb Z)
 \xrightarrow{\sim}
 \{(a,k):a\in\mathbb P^1(\mathbb F_3),\ k\in\mathbb F_3\}
 \tag{GMC-4.11}
\]
qui respecte la fibre de $r$, on associe à chaque direction $\ell$ le projecteur
$P_{\Xi(\ell)}$. En ordonnant les douze directions sous la forme $t=0,\ldots,11$, soit
\[
 F_t=\frac14P_{\Xi(\ell_t)},
 \qquad
 B_{ij}=F_{i+j\bmod12}.
 \tag{GMC-4.12}
\]

\begin{theorem}[Carré qutrit de l'atlas nonadique complet]
\label{p12-83-c22-s3:gmc:thm:qms-12-3}
La matrice $B$ est un \QMS de tailles $(12,3)$, contient les douze projecteurs
étalonnés des quatre contextes qutrit et est semi-classique :
\[
 B=\sum_{t=0}^{11}P_{\sigma_t}\otimes F_t.
 \tag{GMC-4.13}
\]
Elle admet une dilatation simultanée en un \QPM commutatif.
\end{theorem}

\begin{proof}
Chacun des quatre contextes a pour somme $I_3$, donc
$\sum_tF_t=(1/4)\cdot4I_3=I_3$. On applique le
\Cref{p12-83-c22-s3:gmc:lem:qms-ciclico}.
\end{proof}

Le facteur $1/4$ n'est pas décoratif. L'atlas contient quatre résolutions de
l'identité et, sans cette normalisation, la somme serait $4I_3$. L'atlas complet
est donc une mesure à douze résultats, non une cinquième base
orthonormée.

\subsection{Théorèmes opératoriels et critère de fidélité généalogique APP--TPK}

Les trois constructions précédentes ont un statut exact :

\begin{center}
\begin{tabularx}{.96\textwidth}{@{}L{.20\textwidth}L{.21\textwidth}Y@{}}
\toprule
Objet & Théorème clos & Structure HMT non encore transportée\\
\midrule
$A^\Sigma$, $(9,3)$ & feuillet somme, positivité, normalisation et
semi-classicité & feuillet produit, quotient et routes\\
$A^{(9)}$, $(9,3)$ & trois contextes, non-commutativité et dilatation commune &
quatrième contexte, incidence de Hensel, gnomon et mémoire TPK\\
$B$, $(12,3)$ & atlas complet $12\to4\times3$ et dilatation commune &
équivariance de l'étalonnage, dynamique et monodromie\\
\bottomrule
\end{tabularx}
\end{center}

Aucun des deux derniers carrés n'est un \QLS ni un \QPM : leurs cellules sont
des projecteurs multipliés par $1/3$ ou $1/4$, non des projections. Ce sont bien des carrés
magiques quantiques au sens de la définition orthodoxe et ils possèdent bien une dilatation
simultanée. Cette formulation résout l'ambiguïté terminologique et conserve
la différence mathématique entre effet et projection.

\begin{workbox}
L'objectif suivant n'est pas de prouver de nouveau qu'il existe un \QMS d'ordre
neuf. Il s'agit d'en construire un \emph{fidèle à APP--TPK} : outre \textup{(GMC-4.3)}, il doit
porter simultanément $\Sigma$, $\Pi$, quotient, orientation et route, et doit
satisfaire une équation d'équivariance avec l'action de TPK. Un candidat
qui ne satisfait pas cette équation demeurera une ombre valable, mais non l'objet
total.
\end{workbox}


\section[Bigraduation généalogique--matricielle]{Bigraduation
généalogique--matricielle à taille extérieure fixe et trigraduation
\texorpdfstring{$(m,n,s)$}{(m,n,s)} : troncature, compression et dilatation
simultanée}
\label{p12-83-c22-s4:gmc:chap:bigraduacion}

La séparation de De las Cuevas--Netzer entre taille extérieure $n$ et niveau
matriciel $s$ résout une ambiguïté du langage des grilles. HMT exige
une seconde séparation : la profondeur $m$ n'est pas une taille intérieure. Ce
chapitre démontre que les directions $m$ et $s$ sont compatibles et qu'une
famille cohérente à toute profondeur possède une dilatation commune.

\subsection{Système projectif d'histoires et système inductif d'observables}

Soit $(X_m,p_m)_{m\ge0}$ un système d'ensembles finis non vides avec des applications
surjectives
\[
 p_m:X_{m+1}\longrightarrow X_m.
 \tag{GMC-3.1}
\]
Dans HMT, $X_m$ peut être l'ensemble complet des préfixes nonadiques de
longueur $m$ ou un ensemble de survivants pourvu que les restrictions
soient surjectives dans le domaine déclaré. La limite projective
\[
 X_\infty=\varprojlim_m X_m
\]
est l'espace compact des histoires compatibles. Chaque $x\in X_m$ détermine
le cylindre ouvert et fermé
\[
 [x]_m=\{\xi\in X_\infty:\xi_m=x\}.
 \tag{GMC-3.2}
\]

Par dualité, les algèbres diagonales
\[
 D_m=C(X_m)
\]
forment un système inductif au moyen de
\[
 \jmath_m:D_m\longrightarrow D_{m+1},
 \qquad \jmath_m(f)=f\circ p_m.
 \tag{GMC-3.3}
\]
Leur limite uniforme est $C(X_\infty)$ : les fonctions cylindriques forment une
sous-algèbre autoadjointe unitale qui sépare les points et sont denses par
Stone--Weierstrass.

\subsection{La famille à \(2\) degrés}

Pour $s\ge1$, nous définissons l'espace des états matriciels
\[
 \mathcal K_{m,s}=\UCP(D_m,M_s(\mathbb C)).
 \tag{GMC-3.4}
\]
Si $e_x\in D_m$ est la fonction caractéristique de $x\in X_m$, toute
$\Phi\in\mathcal K_{m,s}$ équivaut à une \POVM
\[
 E_x^{(m)}=\Phi(e_x)\succeq0,
 \qquad
 \sum_{x\in X_m}E_x^{(m)}=I_s.
 \tag{GMC-3.5}
\]
La restriction généalogique est
\[
 \rho_m^s:\mathcal K_{m+1,s}\longrightarrow\mathcal K_{m,s},
 \qquad
 \rho_m^s(\Phi)=\Phi\circ\jmath_m.
 \tag{GMC-3.6}
\]
En termes d'effets,
\[
 E_x^{(m)}=\sum_{y\in p_m^{-1}(x)}E_y^{(m+1)}.
 \tag{GMC-3.7}
\]
C'est exactement l'opération « oublier la continuation mais sommer toutes ses
prolongations ». Ce n'est ni une moyenne ni une élimination de masse : elle conserve l'unité.

Dans la direction matricielle, si $d\ge0$,
$\Phi_r\in\mathcal K_{d,s_r}$ et
$V_r:\mathbb C^s\to\mathbb C^{s_r}$ satisfont
$\sum_rV_r^*V_r=I_s$, leur combinaison matriciellement convexe est
\[
 \operatorname{mc}_{\boldsymbol V}(\Phi_1,\ldots,\Phi_k)(f)
 =\sum_{r=1}^kV_r^*\Phi_r(f)V_r
 \in\mathcal K_{d,s}.
 \tag{GMC-3.8}
\]

\begin{theorem}[Commutation de la généalogie et de la convexité matricielle]
\label{p12-83-c22-s4:gmc:thm:conmutacion-ms}
Pour toute famille $\Phi_r\in\mathcal K_{m+1,s_r}$ et toute compression comme dans
\textup{(GMC-3.8)},
\[
 \rho_m^s\!\left(
   \operatorname{mc}_{\boldsymbol V}(\Phi_1,\ldots,\Phi_k)
 \right)
 =
 \operatorname{mc}_{\boldsymbol V}\!\left(
   \rho_m^{s_1}(\Phi_1),\ldots,\rho_m^{s_k}(\Phi_k)
 \right).
 \tag{GMC-3.9}
\]
\end{theorem}

\begin{proof}
Pour $f\in D_m$, les deux membres valent
\[
 \sum_rV_r^*\Phi_r(\jmath_m f)V_r.
\]
L'égalité est fonctorielle, non numérique : elle vaut simultanément pour tous les
niveaux matriciels et pour n'importe quel système de troncatures.
\end{proof}

\begin{bridgebox}
L'équation \textup{(GMC-3.9)} est le noyau de la bigraduation $(m,s)$. Raffiner
l'histoire puis comprimer produit exactement la même ombre que
comprimer chaque lecteur puis oublier la continuation. La profondeur et le
niveau matriciel sont indépendants, mais forment un carré commutatif.
\end{bridgebox}

\begin{figure}[H]
\centering
\begin{tikzpicture}[node distance=2.5cm and 3.7cm,>=Latex]
\node[draw,rounded corners,fill=hmtlightblue,minimum width=3.3cm,
 minimum height=1cm] (a) {$\mathcal K_{m+1,s'}$};
\node[draw,rounded corners,fill=hmtlightblue,right=of a,minimum width=3.3cm,
 minimum height=1cm] (b) {$\mathcal K_{m+1,s}$};
\node[draw,rounded corners,fill=hmtlightgold,below=of a,minimum width=3.3cm,
 minimum height=1cm] (c) {$\mathcal K_{m,s'}$};
\node[draw,rounded corners,fill=hmtlightgold,right=of c,minimum width=3.3cm,
 minimum height=1cm] (d) {$\mathcal K_{m,s}$};
\draw[->,thick] (a)--node[above]{compression} (b);
\draw[->,thick] (a)--node[left]{troncature} (c);
\draw[->,thick] (b)--node[right]{troncature} (d);
\draw[->,thick] (c)--node[below]{compression} (d);
\node[hmtviolet] at ($(a)!0.5!(d)$) {commute};
\end{tikzpicture}
\caption{Les deux directions indépendantes de la bigraduation. Le symbole
$s'$ peut représenter une somme directe de plusieurs niveaux internes.}
\label{p12-83-c22-s4:gmc:fig:cuadrado-ms}
\end{figure}

\subsection{Dilatation uniforme compatible avec toutes les profondeurs du système projectif}

Une famille $(\Phi_m)_{m\ge0}$ est compatible si
\[
 \Phi_m=\Phi_{m+1}\circ\jmath_m.
 \tag{GMC-3.10}
\]
Cela revient à exiger la loi d'agrégation \textup{(GMC-3.7)} pour tous les préfixes.

\begin{theorem}[Dilatation simultanée généalogique, corridor diagonal]
\label{p12-83-c22-s4:gmc:thm:dilatacion-diagonal}
Toute famille compatible
$\Phi_m\in\UCP(C(X_m),M_s)$ détermine une unique application
\[
 \Phi_\infty:C(X_\infty)\longrightarrow M_s
 \tag{GMC-3.11}
\]
qui coïncide avec $\Phi_m$ sur chaque algèbre cylindrique. Il existe un espace de
Hilbert $\mathcal K$, une isométrie $V:\mathbb C^s\to\mathcal K$, une
représentation $\pi:C(X_\infty)\to B(\mathcal K)$ et sa mesure spectrale
projective associée $P$ sur $X_\infty$ tels que
\[
 \Phi_\infty(f)=V^*\left(\int_{X_\infty}f(\xi)\,dP(\xi)\right)V.
 \tag{GMC-3.12}
\]
En particulier, pour tout $m$ et $x\in X_m$,
\[
 E_x^{(m)}=V^*P([x]_m)V.
 \tag{GMC-3.13}
\]
Une \PVM commune dilate simultanément toutes les \POVM finies de la
généalogie. Si la dilatation de Stinespring est choisie minimale, le triplet
$(\pi,\mathcal K,V)$ est unique à équivalence unitaire près ; une fois
$\pi$ fixée, sa mesure spectrale $P$ est unique.
\end{theorem}

\begin{proof}
La compatibilité définit une application unitale positive sur la réunion des
algèbres cylindriques. Puisque le domaine est commutatif, la positivité implique la
positivité complète, et l'application est contractante. Par densité, elle
s'étend de manière unique à $C(X_\infty)$. Le théorème de Stinespring produit
$\Phi_\infty(f)=V^*\pi(f)V$ et garantit que sa réalisation minimale est unique
à équivalence unitaire près. Toute représentation d'une algèbre commutative
est l'intégrale par rapport à une unique \PVM $P$ associée à cette représentation.
Appliquer la formule à la fonction
caractéristique du cylindre $[x]_m$ donne \textup{(GMC-3.13)}.
\end{proof}

\begin{corollary}[Réalisation ambiante fidèle]
La représentation de \Cref{p12-83-c22-s4:gmc:thm:dilatacion-diagonal} peut être choisie fidèle sur
$C(X_\infty)$ sans modifier aucune compression finie.
\end{corollary}

\begin{proof}
Si $(\pi,\mathcal K,V)$ est une dilatation et
$\pi_f:C(X_\infty)\to B(\mathcal K_f)$ une représentation fidèle, on prend
$\pi\oplus\pi_f$ et l'isométrie $V\oplus0$. La compression demeure
$\Phi_\infty$ et la représentation ambiante n'identifie plus deux fonctions
cylindriques distinctes.
\end{proof}

La dernière construction type la phrase « conserver la généalogie ». La
dilatation minimale conserve exactement ce que voit le lecteur ; la dilatation
ambiante fidèle conserve en outre toutes les distinctions de l'algèbre avant
compression. La fidélité d'une dynamique TPK exige cependant également
d'entrelacer ses flèches et ne se déduit pas de la seule fidélité de la
représentation diagonale.

\subsection{Tour UHF non commutative et compatibilité matricielle}

La tour UHF nonadique, sa trace et ses fermetures factorielles sont démontrées dans
\cref{chap:integracion-vn64-operatorial}. Elles sont prises ici comme structure de
départ explicite pour démontrer un résultat distinct : le prolongement
simultané d'une famille compatible d'applications UCP et sa dilatation
commune.

Soit
\[
 \mathcal A_m=M_{9^m}(\mathbb C),
 \qquad \iota_m(a)=a\otimes I_9,
 \qquad
 \mathcal A_\infty=\varinjlim(\mathcal A_m,\iota_m).
 \tag{GMC-3.14}
\]

\begin{theorem}[Dilatation simultanée sur la limite UHF]
\label{p12-83-c22-s4:gmc:thm:dilatacion-uhf}
Si $\Psi_m:\mathcal A_m\to M_s$ sont des applications UCP et
\[
 \Psi_m=\Psi_{m+1}\circ\iota_m,
 \tag{GMC-3.15}
\]
il existe une unique application UCP
$\Psi_\infty:\mathcal A_\infty\to M_s$ qui les prolonge. Il existe une
représentation $\pi:\mathcal A_\infty\to B(\mathcal K)$ et une isométrie
$V:\mathbb C^s\to\mathcal K$ telles que
\[
 \Psi_m(a)=V^*\pi(\iota_{m,\infty}(a))V
 \quad\text{pour tout }m.
 \tag{GMC-3.16}
\]
La représentation peut être rendue fidèle par somme directe avec une représentation
fidèle de la limite UHF.
\end{theorem}

\begin{proof}
La compatibilité définit $\Psi_\infty$ dans la limite inductive algébrique. Les
applications UCP sont complètement contractantes, de sorte qu'elle s'étend à la
complétion $C^*$. Stinespring donne \textup{(GMC-3.16)} ; la somme directe démontre la
dernière affirmation.
\end{proof}

En restreignant \Cref{p12-83-c22-s4:gmc:thm:dilatacion-uhf} aux diagonales canoniques, on retrouve
le cas nonadique complet $X_m=(\mathbb Z/9\mathbb Z)^m$ du
\Cref{p12-83-c22-s4:gmc:thm:dilatacion-diagonal} ; cette restriction ne couvre pas à elle seule un
système arbitraire de survivants. Dans le système UHF nonadique canonique de
\textup{(GMC-3.14)}, la représentation GNS de sa trace unique possède une fermeture faible
isomorphe à $\Rhyper$. Ainsi, la bigraduation n'est pas un appendice combinatoire : elle unit
les cylindres d'histoires, les cartes matricielles finies et le continu
opératoriel de la réalisation de von Neumann.

\subsection{Trigraduation et niveaux de simultanéité}

Lorsqu'une grille apparaît, on incorpore la taille extérieure $n$ :
\[
 \mathfrak M_{m,s}^{(n)}.
 \tag{GMC-3.17}
\]
La structure est trigraduée si $n$ varie et bigraduée par $(m,s)$ si l'on fixe
$n$. « Simultané » peut signifier trois exigences croissantes :

\begin{enumerate}
 \item \emph{par lignes et colonnes} : une dilatation commune des $2n$ \POVM
 d'un carré magique ;
 \item \emph{par profondeurs} : une dilatation commune de tous les niveaux
 $m$, démontrée dans les Théorèmes~\ref{p12-83-c22-s4:gmc:thm:dilatacion-diagonal}
 et~\ref{p12-83-c22-s4:gmc:thm:dilatacion-uhf} ;
 \item \emph{par généalogie dynamique} : en outre, un entrelaceur préservant
 route, monodromie et composition TPK.
\end{enumerate}

Les deux premières possèdent des théorèmes exacts dans leurs domaines. La troisième est la
notion plus forte que HMT apporte au cadre de De las Cuevas--Netzer. La
construction du système d'opérateurs commun et son équation d'équivariance sont
démontrées dans \Cref{sec:cpe-antecedente-ucp} ; la dilatation statique constitue
seulement l'un de leurs antécédents.

\begin{controlbox}
Le théorème n'affirme pas que toute collection disjointe d'effets à différentes
profondeurs possède une dilatation commune. L'hypothèse décisive est la
compatibilité \textup{(GMC-3.10)} ou \textup{(GMC-3.15)}. Dans HMT, cette compatibilité est la
forme opératorielle de « chaque préfixe est exactement la restriction de ses
continuations ».
\end{controlbox}


\section{Cônes matriciels, images et extensions complètement positives}
\label{p12-83-c22-s5:gmc:chap:beyond}

Le programme récent de De les Coves, van der Eyden et Netzer étend les
systèmes d'opérateurs à des systèmes coniques définis sur des catégories plus
générales \citep{delescoves2026beyond}. Ce cadre ne transforme pas n'importe quelle
catégorie en théorie quantique : il exige de choisir des cônes, des dualités, des compressions et
des applications complètement positives appropriés. HMT apporte un terrain particulièrement
riche pour effectuer ce choix.

\subsection{Réalisations dans des systèmes d'opérateurs ordinaires}

Les familles suivantes s'inscrivent entièrement dans le formalisme standard :

\begin{itemize}
 \item $\mathcal S_d=\Span\{I,P_{\Sigma,d},P_{\Pi,d}\}$ et leurs images
 matricielles ;
 \item les algèbres diagonales $C(X_m)$ et la tour UHF
 $M_{9^m}\hookrightarrow M_{9^{m+1}}$ ;
 \item les lecteurs UCP, \POVM, \PVM, les canaux CPTP et les dilatations de
 Stinespring ;
 \item les familles bigraduées $\mathcal K_{m,s}$.
\end{itemize}

Ici, un formalisme catégorique plus large n'est pas nécessaire pour affirmer
la positivité, la dilatation ou la convexité matricielle. En effet, la classification
\textup{(GMC-6.9)} montre qu'une partie substantielle d'APP peut se résoudre par
la théorie opératorielle exacte en dimension finie.

\subsection{Extension généalogique des systèmes d'opérateurs pour les routes TPK}

Une route TPK n'est pas seulement un opérateur entre deux espaces : elle porte une origine, une destination,
un feuillet, une orientation, une longueur, une retenue et une loi de composition. Si l'on oublie
ces étiquettes et ne conserve que la matrice représentant la route,
des flèches distinctes peuvent se confondre. Pour incorporer le type complet, on propose
une catégorie $\mathsf C_{\rm HMT}$ dont les objets sont des profils de registre et
dont les flèches sont des transports admissibles. À chaque type $c$ serait associé un
cône de lecteurs positifs $A(c)$, et à chaque transformation une action
compatible sur ces cônes.

Dans le langage de \emph{Beyond Operator Systems}, l'objectif est de construire un
$S$-système : une extension fonctorielle du produit tensoriel par l'objet de base
à une famille de cônes, avec des applications CP définies comme des transformations
naturelles niveau par niveau. La théorie fournit des systèmes minimaux et maximaux,
la dualité, des théorèmes de Choi--Kraus généralisés, des représentations comme
intersections de $S$-èdres et une catégorie $*$-autonome
\citep{delescoves2026beyond}.

\begin{workbox}
Le $S$-système intrinsèque de TPK reste à construire. Il manque, au minimum : le
foncteur de cônes, l'objet dualisant, les évaluations et coévaluations, et la
preuve que la composition TPK respecte la positivité dans tous les types.
Le fait que TPK soit monoïdale dans sa réalisation finie ne démontre pas à lui
seul qu'elle soit $*$-autonome.
\end{workbox}

\subsection{Cônes généalogiques HMT et extension de la théorie conique}

Les exemples classiques partent d'un cône de base et étudient ses systèmes minimal et
maximal. HMT suggère un cône de base doté d'une mémoire stratifiée :
\[
 \mathcal C_m^{\rm HMT}
 =\operatorname{cone}\{\text{lecteurs positifs d'histoires à la profondeur }m\}.
 \tag{GMC-8.1}
\]
Les troncatures induisent des applications positives
$\mathcal C_{m+1}^{\rm HMT}\to\mathcal C_m^{\rm HMT}$, tandis que les
compressions engendrent les niveaux matriciels. Le nouveau phénomène à étudier est
de savoir si les systèmes minimal et maximal restent séparés lors du passage à la limite en
$m$, et quelle part de cette séparation mesure une perte de généalogie plutôt qu'une simple
non-commutativité.

Cela propose trois invariants pour le programme général :

\begin{enumerate}
 \item \emph{profondeur de séparation} : le plus petit $m$ pour lequel les systèmes
 minimal et maximal diffèrent ;
 \item \emph{niveau matriciel de détection} : le plus petit $s$ qui manifeste cette
 différence ;
 \item \emph{défaut de mémoire} : l'information minimale de route qu'il faut
 ajouter pour relever un point du système maximal au minimal.
\end{enumerate}

Les deux premiers raffinent la séparation $(m,s)$ ; le troisième est proprement
HMT. Ce sont des définitions proposées, qui doivent d'abord être étudiées dans
$\mathfrak C_2^{\rm APP}$, où \textup{(GMC-6.9)} permet des calculs exacts.

\subsection{Construction explicite de cubes magiques quantiques et compatibilité de leurs marges}

L'article sur les carrés latins propose de développer des cubes magiques quantiques.
Le lemme cyclique admet une extension immédiate.

\begin{definition}[Cube magique quantique]
Une famille
$C=(C_{ijk})_{0\le i,j,k<n}$ d'effets dans $M_s$ est un cube magique quantique
si la somme de toute ligne parallèle à l'un quelconque des trois axes vaut $I_s$.
\end{definition}

\begin{theorem}[Cube cyclique d'une \POVM]
\label{p12-83-c22-s5:gmc:thm:cubo-ciclico}
Pour toute \POVM $(E_t)_{t\in\mathbb Z/n\mathbb Z}$,
\[
 C_{ijk}=E_{i+j+k\bmod n}
 \tag{GMC-8.2}
\]
est un cube magique quantique de tailles extérieure $n$ et intérieure $s$.
\end{theorem}

\begin{proof}
Deux coordonnées étant fixées, lorsque la troisième varie, l'indice
$i+j+k$ parcourt une fois $\mathbb Z/n\mathbb Z$. La somme de chaque ligne vaut donc
$\sum_tE_t=I_s$. La positivité est héritée de la \POVM.
\end{proof}

En appliquant le théorème à $E^{(9)}$ de \textup{(GMC-4.6)} et à $F$ de
\textup{(GMC-4.12)}, on obtient des cubes qutrit explicites d'ordres neuf et douze. Le
premier peut consacrer le troisième axe à une coordonnée de phase ; le second utilise
l'atlas nonadique complet. Pour transformer cet axe en profondeur HMT réelle,
les couches doivent en outre satisfaire
\[
 \rho_m(C^{m+1}_{ijk})=C^m_{ijk},
 \tag{GMC-8.3}
\]
et entrelacer TPK. La construction cyclique clôt le cube opératoriel ; la
fidélité généalogique demeure une condition supplémentaire.

\subsection{Combinaisons magiques d'applications CP}

Une autre direction ouverte de la théorie consiste à transposer les conditions
magiques aux applications complètement positives par Choi--Jamiołkowski. Si
$J(\Phi_{ij})\succeq0$ sont les matrices de Choi d'une grille d'applications CP,
les conditions
\[
 \sum_jJ(\Phi_{ij})=J(\id),
 \qquad
 \sum_iJ(\Phi_{ij})=J(\id)
 \tag{GMC-8.4}
\]
sont une version linéaire des normalisations magiques. HMT apporte des canaux de
lecture, de mise à jour et de transport ; le problème est de les décomposer en une
grille qui préserve à la fois la normalisation, le feuillet et la route. Le formalisme
conique généralisé peut traiter des codomaines qui ne sont pas simplement des cônes
PSD matriciels.

\subsection{Critère d'universalité dans le domaine généalogique}

Le cadre de Gonda, Reinhart, Stengele et De les Coves définit un simulateur dans
une catégorie \emph{cible--contexte} par
\[
 s:P\otimes C\longrightarrow T\otimes C,
 \tag{GMC-8.5}
\]
avec un compilateur fonctionnel $s_T:P\to T$ et une réduction de contexte. L'
universalité est l'existence d'une réduction laxe au simulateur trivial,
non la simple richesse ou récurrence du système
\citep{gonda2024universality}.

Une instance HMT peut prendre :

\begin{itemize}
 \item $P$ : préfixes, programmes de route ou germes ;
 \item $T$ : observables mathématiques et physico-métrologiques que l'on cherche à
 réaliser ;
 \item $C$ : jauge, feuillet, conditions aux limites et appareil de lecture ;
\item $s_T$ : transformation d'une généalogie en une sortie ;
 \item $s_C$ : transport du contexte et de la mémoire résiduelle.
\end{itemize}

Pour démontrer l'universalité, il faudrait construire ces données et la réduction
laxe. Pour réfuter une proposition, un monotone du préordre contextuel suffit.
Le théorème d'impossibilité du cadre affirme, en particulier, que si l'image
fonctionnelle du compilateur a un supremum strictement inférieur à l'espace des
cibles pour un certain monotone compatible, le simulateur n'est pas universel.

\begin{bridgebox}
HMT suggère une version plus forte : \emph{l'universalité généalogiquement
fidèle}. Outre reproduire le comportement de la cible, la réduction doit
se relever en un foncteur entre catégories de routes et préserver la composition,
le retour et le registre chronologique enrichi. Deux simulateurs comportementalement équivalents peuvent cesser
d'être équivalents lorsque l'on exige la conservation de la provenance. Cette distinction
offre une instance nouvelle et non triviale au programme catégorique.
\end{bridgebox}

L'unification métrologique démontrée par HMT et l'universalité catégorique
sont des affirmations distinctes. La première dérive des lecteurs et des relations à l'intérieur
du système ; la seconde exige de pouvoir simuler une classe déclarée de cibles et de
contextes. Les typer séparément les renforce toutes deux : cela évite de demander à la métrologie
une propriété dont elle n'a pas besoin et transforme l'universalité en un théorème
réfutable.

\subsection{Contextualité et non-semi-classicité}

Les carrés cycliques de \Cref{p12-83-c22-s3:gmc:chap:qms} sont semi-classiques. Ils ne
résolvent donc pas la direction ouverte concernant la contextualité ou les QMS non
semi-classiques. APP contient cependant la paire non commutative $P_\Sigma,P_\Pi$
et l'algèbre \textup{(GMC-6.9)}. La bonne route consiste à construire des effets de cellule qui
dépendent des deux opérateurs et à prouver qu'ils n'admettent pas la décomposition
\textup{(GMC-4.4)} avec une \POVM sur $S_n$.

La vérification est un problème semi-défini : chercher
$Q_\sigma\succeq0$, $\sum_\sigma Q_\sigma=I$ et
\[
 A_{ij}=\sum_{\sigma:\,\sigma(i)=j}Q_\sigma.
 \tag{GMC-8.6}
\]
La faisabilité certifie la semi-classicité ; l'infaisabilité doit s'accompagner d'un
témoin dual. Une matrice non commutative ou un graphique visuel ne remplacent pas
ce certificat.


\section{Structure généalogique d'APP--TRIT--TPK au niveau matriciel}
\label{p12-83-c22-s6:gmc:chap:hmt-profundidad}

Ce chapitre isole les structures qui interviennent dans la transduction et
garde leurs types explicites. L'ordre est
délibéré : support, différence somme--produit, mémoire ternaire, catégorie de
routes, réalisation quantique, clôture opératorielle et lecture métrologique. Ce n'est
qu'ensuite que l'on demandera lesquelles de ces structures forment un carré magique ou
un système d'opérateurs.

\subsection{Support commun d'APP et évaluation duale par les feuillets somme et produit}

Soit
\[
 X_9=(\mathbb Z/9\mathbb Z)^2,
 \qquad
 \mathscr D=\{(1,0),(0,1),(-1,0),(0,-1)\}.
\]
Le graphe de Cayley $\operatorname{Cay}(X_9,\mathscr D)$ possède 81 sommets et
quatre arêtes cardinales orientées depuis chaque sommet. L'architecture APP
contient tous les chemins finis de ce graphe. Le sigle est conservé comme
nom technique et n'est pas utilisé ici pour fixer un développement historique. À la longueur
$r$, il y a exactement $81\cdot4^r$ chemins si les retours et les
reculs sont permis. Il n'existe pas une table additive et une autre multiplicative indépendantes :
il existe un support unique et deux évaluations du même point,
\[
 \Sigma(i,j)=\operatorname{dr}(i+j),
 \qquad
 \Pi(i,j)=\operatorname{dr}(ij),
 \qquad
 \operatorname{dr}(a)=1+((a-1)\bmod9).
 \tag{GMC-2.1}
\]
La table de $\Sigma$ est un carré latin classique d'ordre neuf, car
toute translation permute $\mathbb Z/9\mathbb Z$. La table de $\Pi$ ne l'est pas :
seules ses lignes et ses colonnes indexées par les six unités
$\{1,2,4,5,7,8\}$ sont des permutations ; $3,6,9$ replient des classes.

\begin{resultbox}
APP contient déjà une structure latine exacte, mais APP ne s'y réduit pas.
Le feuillet somme est latin ; le feuillet produit enregistre le défaut de latinité de l'
anneau local ; les chemins conservent l'ordre de consultation, et le relèvement
reste--quotient conserve ce que la réduction modulaire masquerait.
\end{resultbox}

Pour chaque cellule, on conserve le reste et le quotient :
\[
 i+j=R^+_{ij}+9q^+_{ij},
 \qquad
 ij=R^\times_{ij}+9q^\times_{ij}.
 \tag{GMC-2.2}
\]
Les totaux certifiés sont
\[
 \sum R^+=405,
 \quad \sum R^\times=459,
 \quad \sum q^+=45,
 \quad \sum q^\times=174,
\]
et donc
\[
 2025-810=1215=(459-405)+9(174-45)=54+9\cdot129.
 \tag{GMC-2.3}
\]
La différence visible $54$ et la différence de quotient $129$ ne sont pas deux
noms pour la même donnée. L'une réside dans les restes publiés ; l'autre, dans
la mémoire nécessaire pour recomposer la valeur entière.

\subsection{Hiérarchie spectrale \(4\to2\to6\to54\) du bloc central d'APP}

Le bloc central d'APP prend la forme
\[
 B_c=\begin{pmatrix}7&2\\2&7\end{pmatrix}
     =7I+2R
     =9P_++5P_-,
 \qquad
 R=\begin{pmatrix}0&1\\1&0\end{pmatrix},
 \quad P_\pm=\frac{I\pm R}{2}.
 \tag{GMC-2.4}
\]
On obtient ainsi exactement la hiérarchie typée induite par la
décomposition spectrale précédente :
\[
 \underbrace{9-5}_{4}
 \ ;\ 
 \underbrace{2}_{\text{couplage hors diagonale}}
 \ ;\ 
 \underbrace{51-45}_{6}
 \ ;\ 
 \underbrace{9\cdot6}_{54}.
 \tag{GMC-2.5}
\]
Le premier nombre est le saut spectral primordial ; le deuxième, le couplage
du bloc ; le troisième, l'excès agrégé du secteur non unitaire ; le quatrième,
son transport à l'échelle globale de la table. Ils ne doivent pas être confondus du fait de leur
parenté arithmétique.

Comme $\det B_c=45$, la normalisation a pour déterminant un et appartient au
groupe hyperbolique connexe :
\[
 \frac{B_c}{\sqrt{45}}
 =\exp(\eta R),
 \qquad
 \eta=\frac12\log\frac95,
 \qquad
 \frac59=e^{-2\eta}.
 \tag{GMC-2.6}
\]
La même paire spectrale admet ainsi une lecture contractante et une lecture
lorentzienne. C'est la connexion précise avec la convexité et l'intérieur
hyperbolique que le pont doit conserver : la contraction $5/9$ est l'ombre
de la séparation des deux idempotents $P_\pm$, non une analogie verbale avec
la courbure.

\subsection{TRIT : état visible, orientation, retenue et mémoire généalogique}

TRIT est la signature orientée $\{-1,0,+1\}$ obtenue à partir de $\mathbb F_3$ au moyen de
$0\mapsto0$, $1\mapsto+1$, $2\mapsto-1$. Pour typer son relèvement local,
nous définissons l'alphabet relevé
\[
 \mathcal L_1=\{\nu=r+3c:r,c\in\{-1,0,+1\}\}
 =\{-4,-3,\ldots,4\},
 \qquad
 \widehat\tau(\nu)=(r,c).
 \tag{GMC-2.7}
\]
La décomposition $\nu=r+3c$ est unique dans $\mathcal L_1$ ; $r$ est le reste
visible et $c$ la retenue. Le symbole $\nu$ est une amplitude relevée et ne
s'identifie pas, sans données supplémentaires, à la somme de deux trits visibles. En particulier,
les amplitudes $3,0,-3$ publient le même reste zéro et donnent trois antécédents
distincts,
\[
 0^+=(0,+1),\qquad 0^0=(0,0),\qquad0^-=(0,-1).
 \tag{GMC-2.8}
\]
L'état complet peut contenir
\[
 x=(t,\mu,o,h,c,\sigma,\lambda,g),
 \tag{GMC-2.9}
\]
où $t$ est la signature visible, $\mu$ le feuillet, $o$ l'orientation, $h$ l'
extension, $c$ la retenue, $\sigma$ la frontière, $\lambda$ la route et $g$ la
phase nonadique. La projection $x\mapsto t$ élimine toutes les autres
coordonnées. C'est le point structurel que HMT apporte à la théorie des
dilatations : deux effets observables égaux peuvent posséder des généalogies
distinctes.

Les trois signatures admettent les algèbres réelles
\[
 \mathbb A_\tau=\mathbb R[J]/(J^2+\bar\tau),
\]
de type elliptique ($J^2=-I$), parabolique ($J^2=0$) et hyperbolique
($J^2=+I$). Dans le dernier cas, $P_\pm=(I\pm J)/2$ séparent deux directions
propres. Le TRIT n'est pas seulement le signe : c'est l'observable minimal d'un transport
qui conserve étiquette, orientation et mémoire.

\subsection{TPK : catégorie de routes avec mémoire}

TPK, \emph{Topological Phase Kernel}, vit sur la catégorie libre des
chemins d'APP et sur sa lecture bivariée somme--produit. Deux chemins qui
aboutissent à la même cellule peuvent différer par le mot cardinal, l'enroulement,
le feuillet, le quotient et l'holonomie. TPK conserve cette différence et compose les routes de
manière associative. Le cocycle de retenue
\[
 \kappa_9(a,b)=
 \frac{\operatorname{dr}(a)+\operatorname{dr}(b)
       -\operatorname{dr}(a+b)}9
\]
satisfait
\[
 \kappa_9(a,b)+\kappa_9(a+b,c)
 =\kappa_9(b,c)+\kappa_9(a,b+c),
 \tag{GMC-2.10}
\]
et garantit que la récupération ne dépend pas du parenthésage d'une
concaténation.

Le quotient phase--longueur possède des objets $(b,r)\in\mathbb Z/12\mathbb Z
\times\{0,\ldots,8\}$ et des flèches
\[
 (b,r)\xrightarrow{\ell}
 \left(b+\left\lfloor\frac{r+\ell}{9}\right\rfloor,
       r+\ell\bmod9\right).
 \tag{GMC-2.11}
\]
Après neuf itérations unitaires ($\ell=1$), $r$ revient et une phase
dodécaphasique avance ; après $108=9\cdot12$ itérations unitaires, le quotient revient.
L'état enrichi du TPK peut ne pas
être revenu, car il conserve une mémoire supplémentaire. Le quotient prédictif
minimal des observables publiés possède 36 états et des fibres de cardinal
trois. Ce fait rend visible la différence entre égalité de sortie et
équivalence d'histoire.

\begin{intuitionbox}
Un carré magique organise simultanément les lignes et les colonnes. TPK organise
simultanément la position, le feuillet, la phase et la composition. La traduction correcte ne
consiste pas à appeler « quantique » toute table HMT, mais à construire un lecteur
positif qui préserve les identités d'APP et à déclarer, coordonnée par
coordonnée, quelle mémoire il conserve et quelle mémoire il comprime.
\end{intuitionbox}

\subsection{Réalisation hilbertienne et structure quantique}

La clôture quantique finie de HMT construit, pour chaque profondeur
$d$, l'espace
\[
 \mathcal H_d=\ell^2((\mathbb Z/9\mathbb Z)^d),
 \qquad
 \mathcal A_d=B(\mathcal H_d)\cong M_{9^d}(\mathbb C).
 \tag{GMC-2.12}
\]
Pour $u,v\in(\mathbb Z/9\mathbb Z)^d$, les translations et les phases nonadiques
agissent par
\[
 X_u|x\rangle=|x+u\rangle,
 \qquad
 Z_v|x\rangle=\omega^{v\cdot x}|x\rangle,
 \qquad \omega=e^{2\pi i/9},
\]
et satisfont
\[
 Z_vX_u=\omega^{u\cdot v}X_uZ_v.
\]
La famille complète --- de manière équivalente, les $d$ paires coordonnées
$X_j,Z_j$ --- engendre $\mathcal A_d$ et agit irréductiblement. Pour le caractère
central primitif fixé par $\omega$, la représentation finie de Heisenberg
est déterminée, à équivalence unitaire près. Dans cette réalisation, on
construit des états, des observables autoadjoints, des mesures
spectrales, des probabilités de Born, une dynamique unitaire et des canaux de Kraus ; la
composition des registres se transporte de manière monoïdale au produit tensoriel.
La représentation hilbertienne finie n'est pas présupposée : elle est construite
explicitement et transporte la généalogie du système.

Dans le quotient qutrit apparaissent quatre sous-algèbres abéliennes maximales du
système de Weyl de dimension trois. Leurs quatre bases spectrales sont
mutuellement non biaisées :
\[
 \Tr(P_{a,k}P_{b,\ell})=
 \begin{cases}
  \delta_{k\ell},&a=b,\\[2pt]
  1/3,&a\ne b,
 \end{cases}
 \qquad a,b\in\mathbb P^1(\mathbb F_3).
 \tag{GMC-2.13}
\]
Les douze projecteurs minimaux $P_{a,k}$ constituent l'atlas opératoriel qui sera
plus loin étalonné avec les douze directions de
$\mathbb P^1(\mathbb Z/9\mathbb Z)$.

\subsection{De la tour nonadique au continu opératoriel}

L'inclusion unitale, la branche marquée, la limite UHF et les fermetures de types
\(I_\infty\) et \(II_1\) sont construites dans
\cref{chap:integracion-vn64-operatorial}. Leurs formules sont des prémisses du pont
matriciel et déterminent la bigraduation.

La tour opératorielle prolonge les cartes finies par les
inclusions unitaires préservant la trace normalisée
\[
 \iota_d:M_{9^d}(\mathbb C)\longrightarrow M_{9^{d+1}}(\mathbb C),
 \qquad a\longmapsto a\otimes I_9.
 \tag{GMC-2.14}
\]
En effet, si $\tau_n=\Tr/n$, alors
$\tau_{9^{d+1}}\circ\iota_d=\tau_{9^d}$ ; la trace ordinaire, en revanche, est
multipliée par neuf. 
Sa limite inductive est l'algèbre UHF de type $9^\infty=3^\infty$. La fermeture
dans la représentation GNS de la trace unique est le facteur hyperfini
$\Rhyper$ de type $\mathrm{II}_1$. Les inclusions marquées
$a\mapsto a\otimes p_j$ suivent une branche et conduisent au corridor des opérateurs compacts,
de type $\mathrm I_\infty$. Le résultat distingue donc la conservation
simultanée de toutes les branches de la sélection d'une histoire précise.

Dans le corridor diagonal, les ensembles finis de préfixes forment un système
projectif $X_{m+1}\to X_m$. L'espace des histoires infinies est
$X_\infty=\varprojlim X_m$, tandis que leurs algèbres de fonctions forment le
système inductif opposé. Cette dualité des directions --- projective pour les
histoires, inductive pour les observables --- sera le support exact de la nouvelle
bigraduation.

\subsection{Double projection et métrologie}

HMT n'identifie pas le nombre réel publié à l'état qui l'engendre. Une
famille compatible de cylindres produit une lecture archimédienne ; une autre
projection organise la sortie exceptionnelle et ses symétries. La réalité
numérique apparaît comme l'ombre d'une double projection d'un état plus riche.
Les théorèmes de génération monodromique de $\pi$, $e$, $\varphi$ et $\alpha$,
les lecteurs d'action et de
gravitation, la loi des masses, les secteurs de mélange et de violation CP, et la
généalogie exceptionnelle s'appliquent dans leurs domaines déclarés. Le pont rend
explicites les objets opératoriels qui transportent leurs conclusions.

L'identité de Planck utilise le même bloc $B_c$, son
transport dans $SO^+(1,1)$ et l'identité autoduale de la droite des masses. Le
fait mathématique décisif pour cette traduction est que « observable physique » et
« histoire génératrice » sont des types différents. Une application UCP peut publier
le premier ; la coordonnée $m$ et la catégorie TPK conservent la seconde.

\begin{sourcebox}
Les prémisses du pont sont l'inclusion unitale
$a\mapsto a\otimes I_9$, le bloc $B_c=9P_++5P_-$, l'identité autoduale de
la droite des masses et la périodicité d'ordre \(108\). Chacune est utilisée avec
le domaine et les hypothèses de son théorème ; les conclusions matricielles se
déduisent au moyen des applications écrites dans ce chapitre.
\end{sourcebox}


% El capítulo se distribuye en dos dependencias causales. La unión de los dos
% selectores reproduce todas sus líneas científicas exactamente una vez.
\ifdefined\HMTGMCContinuumMPS
\section{Limites généalogiques, canaux de transfert et mémoire de Stinespring}
\label{p12-83-c22-s7:gmc:chap:continuo-mps}

Le continu HMT et la limite continue des états de produit matriciel (MPS)
répondent à des questions différentes. Le premier prolonge des histoires finies et
les lit ensuite ; la seconde raffine un réseau tensoriel et exige que son canal de
transfert admette des racines compatibles. La comparaison n'est féconde que si l'on
construit le canal qui les relie.

\subsection{Séparation typée de la limite généalogique et de la limite MPS}

Dans le corridor HMT, une histoire infinie est un élément de
$X_\infty=\varprojlim X_m$. Sous les conditions de finitude, de surjectivité,
de compacts emboîtés, de contraction et de couverture cohérente démontrées dans l'
annexe sur le continu, le quotient par la lecture archimédienne reconstruit le
compact réel déclaré. L'état complet et sa valeur réelle demeurent des
objets distincts.

Pour une MPS invariante par translation, un tenseur
$A=\{A^i\in M_D\}_{i=1}^d$ détermine le canal de transfert
\[
 E_A=\sum_i A^i\otimes\overline{A^i}.
 \tag{GMC-7.1}
\]
De las Cuevas, Schuch, Pérez-García et Cirac démontrent qu'une MPS sous forme
irréductible possède une limite continue si et seulement si $E_A$ est infiniment
divisible ; de manière équivalente, il existe un canal projecteur $P$ et un générateur
de Lindblad $L$ tels que
\[
 E_A=Pe^L,
 \qquad P^2=P,
 \qquad PLP=PL.
 \tag{GMC-7.2}
\]
Elle possède une limite continue grossière si une certaine puissance $E_A^r$ est infiniment
divisible \citep{delascuevas2018continuum}.

\begin{controlbox}
L'existence de la limite projective HMT n'implique pas \textup{(GMC-7.2)}. Pour appliquer le
théorème MPS, il faut exhiber un tenseur local uniforme, identifier son canal de
transfert et vérifier la divisibilité infinie. Réciproquement,
\textup{(GMC-7.2)} ne conserve pas à elle seule le feuillet, la retenue, la route ou la monodromie.
\end{controlbox}
\fi

\ifdefined\HMTGMCOperatorialTower
\subsection{Raffinement nonadique par espérance conditionnelle}
\label{p12-83-c22-s7:gmc:sec:refinamiento-condicional}

Dans la tour
$\mathcal A_{m+1}=\mathcal A_m\otimes M_9$, l'espérance conditionnelle
\[
 \mathbb E_m=\id\otimes\tau_9:
 \mathcal A_{m+1}\longrightarrow\mathcal A_m,
 \qquad \tau_9(b)=\frac1{9}\Tr(b),
 \tag{GMC-7.3}
\]
est UCP et satisfait
\[
 \mathbb E_m(a\otimes I_9)=a.
 \tag{GMC-7.4}
\]
Sur la diagonale $D_{m+1}=C(X_m\times\mathbb Z_9)$, elle se réduit à la moyenne
uniforme sur le dernier chiffre. L'évaluation de branche
\[
 \kappa_{m,a}(f)(x)=f(x,a)
 \tag{GMC-7.5}
\]
est une autre rétraction, cette fois multiplicative sur la diagonale. Faire la moyenne de toutes
les branches et sélectionner une branche sont des opérations distinctes, parallèles aux
inclusions $a\otimes I_9$ et $a\otimes p_a$ de la clôture opératorielle.

L'espérance \textup{(GMC-7.3)} possède une dilatation avec un registre explicite. Soient
\[
 W_a\psi=\frac13\,\psi\otimes e_a,
 \qquad a=0,\ldots,8.
 \]
Alors
\[
 \mathbb E_m(X)=\sum_{a=0}^8W_a^*XW_a.
 \tag{GMC-7.6}
\]
Sous la forme de Stinespring,
\[
 \pi(X)=X\otimes I_9,
 \qquad
 V\psi=\frac13\sum_{a=0}^8(\psi\otimes e_a)\otimes e_a,
 \qquad
 \mathbb E_m(X)=V^*\pi(X)V.
 \tag{GMC-7.7}
\]
Le système auxiliaire (\emph{ancilla}) final étiquette exactement le chiffre écarté par la moyenne.

\begin{bridgebox}
La formule \textup{(GMC-7.7)} est une mémoire nonadique canonique au sens de
Stinespring : ce qui se perd en appliquant $\mathbb E_m$ reste orthogonalement
enregistré. Elle ne s'identifie pas encore à la mémoire TPK. Cette identification
requiert un entrelaceur qui transporte aussi la route, le feuillet et la dynamique.
\end{bridgebox}

\subsection{Mémoire de Stinespring et croissance non uniforme}
\label{p12-83-c22-s7:gmc:sec:memoria-no-uniforme}

Le \Cref{p12-83-c22-s4:gmc:thm:dilatacion-uhf} donne une représentation commune, en général
de dimension infinie, pour toute famille compatible. On ne peut exiger en
général un système auxiliaire fini fixe qui purifie tous les niveaux. En effet, la
famille compatible de traces normalisées sur $M_{9^m}$ a pour rang $9^m$ au
niveau $m$, et toute purification de l'état correspondant nécessite
une dimension auxiliaire d'au moins $9^m$. Par conséquent,
\[
 \text{compatibilité UCP à toute profondeur}
 \quad\centernot\Longrightarrow\quad
 \text{mémoire auxiliaire finie uniforme}.
 \tag{GMC-7.8}
\]
La croissance de la mémoire n'est pas un défaut de la construction : elle quantifie la
quantité de généalogie que l'on exige de conserver.

La fidélité dynamique de cette mémoire et sa comparaison avec la limite des
états de produit matriciel sont développées dans
\cref{p12-83-c22-s7:gmc:sec:fidelidad-dinamica}, sans répéter la tour UHF ni la dilatation
de Stinespring démontrées ici.
\fi

\ifdefined\HMTGMCContinuumMPS
\subsection{Équation de fidélité dynamique}
\label{p12-83-c22-s7:gmc:sec:fidelidad-dinamica}

L'espérance préservant la trace \(\id\otimes\tau_9\), sa dilatation de
Stinespring et l'impossibilité générale d'une mémoire auxiliaire finie
uniforme ont été construites dans
\cref{p12-83-c22-s7:gmc:sec:refinamiento-condicional,p12-83-c22-s7:gmc:sec:memoria-no-uniforme}. Sur
cette prémisse opératorielle, et sans abréger ses preuves, on formule maintenant la
condition supplémentaire de fidélité généalogique.

Soit $T_m^{\rm TPK}$ le transport entre registres à la profondeur $m$, et soit
$J_m$ l'encodage de ces registres dans l'espace de dilatation. La
mémoire de Stinespring est TPK-fidèle s'il existe un transport
$\widetilde T_m$ dans l'espace ambiant tel que
\[
 J_{m+1}T_m^{\rm TPK}=\widetilde T_mJ_m
 \tag{GMC-7.9}
\]
et si les compressions observationnelles satisfont la naturalité correspondante.
Si deux histoires distinctes sont identifiées par $J_m$, la dilatation peut être
opératoriellement valide mais non généalogiquement fidèle.

L'équation \textup{(GMC-7.9)} est réfutable : il suffit d'une paire de routes que TPK
distingue et que l'espace ambiant identifie, ou d'une composition dont l'image ne coïncide pas
avec la composition des images.

\subsection{Critère de compatibilité avec les limites d'états de produit matriciel}

Pour appliquer \textup{(GMC-7.2)} au continu HMT, cinq étapes doivent être accomplies :

\begin{enumerate}
 \item choisir une cellule locale uniforme et construire des matrices
 $A^i\in M_D$ ;
 \item démontrer que le regroupement nonadique des cellules entrelace l'évolution
 avec une isométrie de raffinement ;
 \item calculer le canal $E_A$ et sa forme irréductible ;
 \item déterminer si $E_A$ ou une certaine puissance $E_A^r$ est infiniment
 divisible ;
 \item transporter à la limite le feuillet, le cocycle de retenue et la
 monodromie, au lieu de ne conserver que l'état tensoriel.
\end{enumerate}

Si le canal s'avérait unitaire, $E_A=\Ad_U$, il aurait des racines unitaires de tout
ordre par calcul fonctionnel. Mais il faudrait encore démontrer que ce canal est
le canal de transfert associé à un tenseur HMT et que ses racines respectent le raffinement
généalogique. Si une partie projective $P$ apparaissait, la mémoire HMT offrirait
une interprétation concrète des secteurs qui survivent au regroupement
d'échelle.

\subsection{Décompositions tensorielles avec symétrie}

Le travail de De las Cuevas avec Netzer et leurs collaborateurs sur les
décompositions tensorielles approchées et sur les complexes simpliciaux
permet une autre entrée \citep{delascuevas2021approximate,
delascuevas2024simplicial}. APP apporte un complexe fini muni d'une action de
translations ; TPK apporte des étiquettes et des compatibilités sur les routes. Une
décomposition tensorielle HMT devrait déclarer :

\begin{itemize}
 \item quel tenseur encode une cellule ou une transition ;
 \item quel groupe préserve les étiquettes et quelle partie n'agit qu'après
 réduction ;
 \item si la factorisation est exacte, approchée ou seulement asymptotique ;
 \item quel rang reste stable lorsque $m$ augmente.
\end{itemize}

La possibilité bidirectionnelle est concrète. Leur formalisme apporte des critères de
factorisation, d'invariance et de limite ; HMT apporte un système de raffinement avec
des fibres, une retenue et une monodromie qui peut produire des exemples où une
factorisation visible existe mais où aucune factorisation généalogiquement fidèle
ne conserve la route.

\begin{workbox}
Un article ultérieur peut étudier la « divisibilité infinie avec mémoire des
racines » : non seulement si $E=E_k^k$, mais si l'on peut choisir une tour cohérente de
racines $E_{k\ell}^\ell=E_k$ munie des registres TPK. Cette condition est
plus forte que la divisibilité du canal nu et constitue une extension
naturelle du problème MPS à partir de HMT.
\end{workbox}
\fi


\section{Réalisation matricielle complètement positive de la fibration nonadique des observables}
\label{p12-83-c22-s8:gmc:chap:simplectica}

La géométrie symplectique ne fait pas partie des théorèmes publiés sur
les carrés magiques de De las Cuevas--Netzer. Le lien légitime passe par une
voie voisine et standard : opérateurs de Weyl, espace des phases discret et
bases mutuellement non biaisées \citep{gibbons2004phase,planat2007rings}. Dans HMT,
cette voie acquiert un relèvement nonadique et une mémoire de fibre.

\subsection{La forme alternée et la commutation de Weyl}

Sur $V_3=\mathbb F_3^2$, nous définissons
\[
 \langle(q,p),(q',p')\rangle_{\rm sp}=qp'-pq'.
 \tag{GMC-5.1}
\]
Si $X|k\rangle=|k+1\rangle$,
$Z|k\rangle=\omega^k|k\rangle$ et $\omega=e^{2\pi i/3}$, les opérateurs
$W_{q,p}=X^qZ^p$ satisfont, avec ces conventions,
\[
 W_uW_v=\omega^{-\langle u,v\rangle_{\rm sp}}W_vW_u.
 \tag{GMC-5.2}
\]
Par conséquent, deux opérateurs commutent exactement lorsque leurs étiquettes sont
orthogonales pour la forme alternée. En dimension deux, les sous-droites
isotropes maximales sont les quatre points de
\[
 \mathbb P^1(\mathbb F_3)=\{[1:0],[0:1],[1:1],[1:2]\}.
 \tag{GMC-5.3}
\]
Chacune engendre un contexte abélien maximal de $M_3(\mathbb C)$ ; ses
projecteurs spectraux constituent une \PVM. Les contextes associés à des droites
distinctes produisent les quatre bases MUB de \textup{(GMC-2.13)}.

\begin{proposition}[Décomposition qutrit par directions]
\label{p12-83-c22-s8:gmc:prop:descomposicion-masa-qutrit}
Soient $\mathcal D_a$ les quatre sous-algèbres diagonales déterminées par les
droites de \textup{(GMC-5.3)} et $\mathcal D_a^0$ leurs parties de trace nulle. Alors
\[
 M_3(\mathbb C)
 =\mathbb CI\oplus
   \bigoplus_{a\in\mathbb P^1(\mathbb F_3)}\mathcal D_a^0
 \tag{GMC-5.4}
\]
comme somme orthogonale pour le produit de Hilbert--Schmidt. Le compte des
dimensions est $9=1+4\cdot2$.
\end{proposition}

\begin{proof}
Chaque $\mathcal D_a^0$ est de dimension deux. L’absence de biais implique
$\Tr(A^*B)=0$ pour $A\in\mathcal D_a^0$ et
$B\in\mathcal D_b^0$, $a\ne b$. Toutes sont orthogonales à $I$. Les
dimensions totalisent neuf, la dimension de $M_3(\mathbb C)$.
\end{proof}

Cette identité explique pourquoi les douze projections ne sont pas douze pièces
indépendantes arbitraires : ce sont quatre résolutions de l’unité dont les parties
centrées décomposent tout l’espace opératoriel qutrit.

\subsection{Le relèvement \texorpdfstring{$12\to4\times3$}{12 vers 4 fois 3}}

Sur le module nonadique $V_9=(\mathbb Z/9\mathbb Z)^2$, la forme
alternée est conservée
\[
 \langle(a,b),(c,d)\rangle_9=ad-bc\pmod9.
 \tag{GMC-5.5}
\]
La droite projective $\mathbb P^1(\mathbb Z/9\mathbb Z)$ est formée des
vecteurs primitifs, modulo la multiplication par une unité. Elle peut
être représentée comme
\[
 \{[1:t]:t\in\mathbb Z/9\mathbb Z\}
 \ \sqcup\ 
 \{[3u:1]:u\in\mathbb Z/3\mathbb Z\},
 \tag{GMC-5.6}
\]
et possède $9+3=12$ points. La réduction modulo trois donne le diagramme
\[
 \mathbb P^1(\mathbb Z/9\mathbb Z)
 \xrightarrow{\ r\ }
 \mathbb P^1(\mathbb F_3),
 \qquad |r^{-1}(a)|=3.
 \tag{GMC-5.7}
\]

La base $a$ d’un contexte qutrit et le résultat $k$ de sa \PVM
forment également un ensemble $4\times3$. Un \emph{étalonnage} est une bijection
\[
 \Xi:\mathbb P^1(\mathbb Z/9\mathbb Z)
 \longrightarrow
 \{P_{a,k}:a\in\mathbb P^1(\mathbb F_3),\ k\in\mathbb F_3\}
 \tag{GMC-5.8}
\]
qui envoie chaque fibre de $r$ vers les trois résultats du contexte situé au-dessus de
$a$.

\begin{controlbox}
La fibration est canonique ; l’affectation de l’ordre interne de chaque fibre aux
trois résultats $k$ est une jauge. Une bijection n’est pas encore une
équivariance. Pour promouvoir $\Xi$ en entrelaceur, il faut démontrer que les
actions linéaires, projectives et antiunitaires pertinentes sont transportées dans
toutes les fibres. La jauge préliminaire échoue à une vérification antiunitaire dans une
fibre ; cet échec est conservé comme réfutateur, il n’est pas dissimulé.
\end{controlbox}

\subsection{Profondeur, retenue et mémoire sur l’espace des phases discret}

La géométrie de Weyl sur $\mathbb F_3^2$ explique la commutation, les contextes et les
MUB. HMT ajoute trois niveaux qui ne peuvent être retrouvés à partir de la seule forme
symplectique :

\begin{enumerate}
 \item le relèvement d’une direction ternaire vers trois points de
 $\mathbb P^1(\mathbb Z/9\mathbb Z)$ ;
 \item le feuillet, la retenue et l’orientation qui distinguent les états de
 même réduction ;
 \item la route TPK qui transporte ces coordonnées et peut revenir à la
 phase sans revenir à l’état.
\end{enumerate}

L’apport potentiel à la littérature qutrit n’est pas « une autre construction de
quatre MUB », déjà connue. C’est une géométrie de \emph{racines de contexte} : chaque
contexte ternaire possède trois antécédents nonadiques, et ces antécédents
peuvent porter une dynamique et une monodromie. Le carré d’ordre douze du
\Cref{p12-83-c22-s3:gmc:thm:qms-12-3} est la première ombre normalisée de cet atlas.

\subsection{Séparation typée entre le recouvrement qutrit
\texorpdfstring{$1/3$}{1/3} et le lecteur HMT
\texorpdfstring{$1/\pi$}{1/pi}}

Dans une paire de bases MUB qutrit,
\[
 \Tr(P_{a,k}P_{b,\ell})=\frac13.
 \tag{GMC-5.9}
\]
Dans le lecteur HMT défini ici, la relation entre un feuillet réel et une
lecture volumétrique peut prendre le quotient
\[
 \frac{\pi}{\pi^2}=\frac1\pi.
 \tag{GMC-5.10}
\]
Les deux expressions ne sont pas interchangeables : $1/3$ est un recouvrement entre
projecteurs en dimension trois ; $1/\pi$ est un poids scalaire d’un autre lecteur.
La différence n’empêche pas de les coordonner. Une \PVM
$\{P_0,P_1,P_2\}$ étant fixée, nous définissons une mesure à neuf effets :
\[
 F_k^{(\pi)}=\frac3\pi P_k\quad(k=0,1,2),
 \qquad
 F_r^{(\pi)}=\frac{1-3/\pi}{6}I_3\quad(r=3,\ldots,8).
 \tag{GMC-5.11}
\]
C’est une \POVM parce que $3/\pi<1$ et
\[
 \sum_{r=0}^{8}F_r^{(\pi)}
 =\frac3\pi I_3+\left(1-\frac3\pi\right)I_3=I_3.
 \tag{GMC-5.12}
\]
Si $Q$ est un projecteur de rang un d’un contexte MUB distinct,
\[
 \Tr(QF_k^{(\pi)})
 =\frac3\pi\Tr(QP_k)
 =\frac3\pi\cdot\frac13
 =\frac1\pi.
 \tag{GMC-5.13}
\]

\begin{bridgebox}
La formule \textup{(GMC-5.11)} transforme exactement la probabilité MUB $1/3$ en
le lecteur $1/\pi$ au moyen de la jauge $3/\pi$, sans identifier les deux nombres.
Les six effets résiduels conservent l’unité et complètent une \POVM à
neuf résultats, à partir de laquelle on obtient un autre \QMS cyclique $(9,3)$. Cette
construction n’identifie pas encore la provenance géométrique de la masse
résiduelle ; elle fixe le type opératoriel exact qu’elle devra avoir.
\end{bridgebox}

\subsection{Groupe symplectique et équation d’équivariance}

En dimension deux sur $\mathbb F_3$, le groupe symplectique coïncide avec
$SL(2,\mathbb F_3)$ et permute les quatre droites. Son action se relève,
projectivement, par le groupe de Clifford sur les \PVM qutrit. Soit
$G_{\rm HMT}$ un groupe ou groupoïde de transformations de l’atlas nonadique et
$U_g$ l’opérateur unitaire ou antiunitaire candidat dans la carte qutrit. La
condition qui transforme un étalonnage en un pont dynamique est
\[
 \Xi(g\cdot\ell)=U_g\,\Xi(\ell)\,U_g^*,
 \qquad
 g\in G_{\rm HMT},\quad
 \ell\in\mathbb P^1(\mathbb Z/9\mathbb Z).
\tag{GMC-5.14}
\]
Si $g$ est une flèche et non seulement une symétrie globale, l'équation doit inclure
la source et le but et se formule comme une transformation naturelle. La vérification de
\textup{(GMC-5.14)} est l'étape qui distingue un étiquetage heureux d'une
transduction structurelle.

\begin{sourcebox}
Les identités symplectiques et MUB s'appuient sur la littérature relative à la phase finie
\citep{gibbons2004phase,planat2007rings}. Elles ne sont pas attribuées aux articles de
De las Cuevas--Netzer sur les carrés magiques. La fibration nonadique, l'étalonnage
et les lecteurs de route constituent la contribution propre de la
construction HMT étudiée ici.
\end{sourcebox}
